\documentclass[aspectratio=169,11pt]{beamer}
\usepackage[utf8]{inputenc}
\usepackage[T1]{fontenc}
\usepackage[spanish]{babel}
\usepackage{lmodern}
\usepackage{booktabs}
\usepackage{tabularx}
\usepackage{hyperref}
\usepackage{enumitem}

\usetheme{Madrid}
\usecolortheme{default}
\setbeamertemplate{navigation symbols}{}
\setbeamertemplate{footline}[frame number]

\title{IA agéntica para la creación de contenido académico}
\subtitle{Sesión 1 --- Fundamentos agénticos e integridad (EIC)}
\author{Juliho Castillo}
\institute{CADi 40\,h · Escuela de Ingeniería y Ciencias\\
Tecnológico de Monterrey}
\date{Clave sugerida: \texttt{CADI-EIC-IA-AGENTICA-01}}

\begin{document}

% --- 1 Portada ---
\begin{frame}
\titlepage
\end{frame}

% --- 2 Mapa 40h ---
\begin{frame}{Mapa del CADi (40 horas)}
\begin{itemize}
\item \textbf{24\,h en aula} $\rightarrow$ 8 sesiones $\times$ 3\,h
\item \textbf{16\,h fuera de aula} $\rightarrow$ lecturas, práctica y microartefacto
\item Hoy: \textbf{Sesión 1 / 8} --- bases + integridad EIC
\item Producto del día (\textbf{E1}): mapa + nota de integridad + tema tentativo
\end{itemize}
\begin{block}{Diseño del CADi}
El trabajo independiente es parte del diseño, no ``tarea opcional''.
\end{block}
\end{frame}

% --- 3 RA ---
\begin{frame}{Resultados de aprendizaje (sesión)}
Al finalizar, usted podrá:
\begin{enumerate}
\item Explicar \textbf{chat ad hoc} vs.\ \textbf{flujo agéntico} con un ejemplo EIC propio
\item Identificar \textbf{roles} (planificador, redactor, revisor, verificador) y puntos \textbf{HITL}
\item Redactar una \textbf{nota de integridad} (qué automatiza / qué decide la persona)
\item Proponer un \textbf{tema EIC tentativo} para el microartefacto
\end{enumerate}
\end{frame}

% --- 4 Normas ---
\begin{frame}{Normas de taller e integridad}
\begin{itemize}
\item Material propio EIC: capítulo, guía de lab, deck o sección de artículo
\item Supervisión humana obligatoria en hechos, ecuaciones, citas y datos
\item Declarar asistencia de IA cuando publique o imparta el material
\item \textbf{No improvisar normativa:} política de integridad / uso de IA $\rightarrow$ verificar con \textbf{CEDDIE} \emph{(enlace y vigencia por edición)}
\end{itemize}
\end{frame}

% --- 5 Chat vs agente ---
\begin{frame}{De chat a agente}
\begin{tabularx}{\textwidth}{@{}XX@{}}
\toprule
\textbf{Chat ad hoc} & \textbf{Flujo agéntico} \\
\midrule
Una conversación suelta & Roles y restricciones explícitas \\
Contexto frágil & Memoria / brief reutilizable \\
``Pide y copia'' & Herramientas + criterio de parada \\
Difícil de auditar & Decisiones humanas visibles (HITL) \\
\bottomrule
\end{tabularx}
\vspace{1em}
\textbf{Definición operativa:} un agente actúa con \textbf{rol}, \textbf{límites}, \textbf{herramientas} y \textbf{parada} ---bajo aprobación humana.
\end{frame}

% --- 6 Anatomía ---
\begin{frame}{Anatomía de un agente}
\begin{itemize}
\item \textbf{Rol:} qué hace (planificar, redactar, revisar\ldots)
\item \textbf{Restricciones:} nivel, notación, qué \emph{no} inventar
\item \textbf{Memoria / contexto:} brief, glosario, RA del curso
\item \textbf{Herramientas:} plantillas, búsqueda, checklists
\item \textbf{Criterio de parada:} cuándo detenerse y pedir humano
\end{itemize}
\begin{alertblock}{Sin estas piezas}
Solo hay un chat con disfraz de ``agente''.
\end{alertblock}
\end{frame}

% --- 7 Orquestación preview ---
\begin{frame}{Preview: orquestación (sesión 2)}
\begin{center}
\texttt{Usuario $\rightarrow$ Orquestador $\rightarrow$ Roles $\rightarrow$ Herramientas $\rightarrow$ Humano aprueba}
\end{center}
\begin{itemize}
\item Secuencial: plan $\rightarrow$ borrador $\rightarrow$ revisión
\item Crítico-humano: el humano es un ``nodo'' obligatorio
\item Paralelo ligero: dos roles, una fusión supervisada
\end{itemize}
Hoy solo el mapa; el diseño detallado es de la \textbf{sesión 2}.
\end{frame}

% --- 8 HITL ---
\begin{frame}{Human-in-the-loop (HITL)}
Inserte aprobación humana \textbf{antes} de:
\begin{itemize}
\item Afirmar ecuaciones, unidades o resultados numéricos
\item Citar fuentes o datos empíricos
\item Publicar / subir a LMS / compartir con estudiantes
\item Cambiar notación o glosario de un curso
\end{itemize}
\begin{alertblock}{STEM académico}
HITL no es ``opcional'' en contenido académico STEM.
\end{alertblock}
\end{frame}

% --- 9 Roles STEM ---
\begin{frame}{Roles útiles en contenido STEM (EIC)}
\begin{tabularx}{\textwidth}{@{}lX@{}}
\toprule
Rol & Pregunta que responde \\
\midrule
Planificador & ¿Qué estructura y RA? \\
Redactor & ¿Qué texto / sección? \\
Revisor & ¿Qué falla en claridad o tono? \\
Verificador & ¿Ecuaciones, unidades, citas OK? \\
\bottomrule
\end{tabularx}
\vspace{1em}
Combine 2--4 roles; no necesite un ejército el día 1.
\end{frame}

% --- 10 Riesgos ---
\begin{frame}{Riesgos con IA en materiales EIC}
\begin{itemize}
\item \textbf{Alucinación:} hechos, fórmulas o procedimientos inventados
\item \textbf{Citas fantasma:} referencias que no existen
\item \textbf{Sesgo / tono} inadecuado al nivel del curso
\item \textbf{Datos sensibles} pegados al chat sin control
\item \textbf{Autoría opaca:} no se sabe qué decidió la persona
\end{itemize}
Mitigar = restricciones + HITL + registro (agent log).
\end{frame}

% --- 11 Ecuaciones ---
\begin{frame}{Ecuaciones, unidades y citas}
Checklist rápido antes de aceptar una salida:
\begin{itemize}
\item[$\square$] Símbolos coherentes con el glosario del curso
\item[$\square$] Unidades SI (o las del programa) sin mezclar
\item[$\square$] Cada cita existe y dice lo que se afirma
\item[$\square$] Afirmaciones empíricas con fuente o marcadas \texttt{[VERIFICAR]}
\end{itemize}
En duda: \textbf{no publicar}; marcar y verificar con experta/o.
\end{frame}

% --- 12 Transparencia ---
\begin{frame}{Transparencia y agent log (adelanto)}
Documente en pocas líneas:
\begin{itemize}
\item Herramienta / modelo usado
\item Rol del agente y prompt/brief
\item Qué generó la IA
\item Qué cambió o rechazó la persona
\item Qué quedó pendiente de verificar
\end{itemize}
Sesiones \textbf{7--8} profundizan; empiece el hábito \textbf{hoy}.
\end{frame}

% --- 13 Demo ---
\begin{frame}{Demo --- idea vaga vs.\ flujo con HITL}
\textbf{Sin flujo:} ``Hazme el capítulo de termodinámica'' $\rightarrow$ texto largo, sin RA, citas dudosas.

\textbf{Con flujo:}
\begin{enumerate}
\item Brief (audiencia, alcance, restricciones)
\item Agente propone outline
\item \textbf{Humano aprueba} nodos y notación
\item Redacción por secciones cortas
\item Verificación explícita (ecuaciones / citas)
\end{enumerate}
\end{frame}

% --- 14 Anti-patrones ---
\begin{frame}{Anti-patrones a evitar}
\begin{itemize}
\item Prompts eternos sin brief ni criterio de parada
\item Cero restricciones (``sé creativo'' en material técnico)
\item Aceptar la primera salida sin editar
\item Inventar bibliografía ``para que se vea serio''
\item Pegar datos de estudiantes o propiedad sensible en el chat
\end{itemize}
Si ocurre: pare, documente y rediseñe el punto HITL.
\end{frame}

% --- 15 Actividad ---
\begin{frame}{Actividad de taller ($\sim$50 min)}
\textbf{Entregable E1} (un solo archivo):
\begin{enumerate}
\item \textbf{Mapa} Usuario $\rightarrow$ Orquestador $\rightarrow$ Roles $\rightarrow$ Herramientas $\rightarrow$ Aprobación
\item \textbf{Nota de integridad} (5--8 líneas): qué automatiza / qué verifica la persona
\item \textbf{Tema EIC tentativo:} audiencia + alcance (una frase cada uno)
\end{enumerate}
Nombre sugerido: \texttt{Apellido\_Sesion1\_integridad-tema}
\end{frame}

% --- 16 Rúbrica ---
\begin{frame}{Rúbrica rápida (E1)}
\begin{tabularx}{\textwidth}{@{}lXX@{}}
\toprule
Criterio & Bien & A mejorar \\
\midrule
Alineación EIC & Audiencia / notación / objetivos claros & Genérico, sin curso \\
HITL & Decisiones y \texttt{[VERIFICAR]} visibles & Aceptación pasiva \\
Reutilización & Sirve a más de una unidad & Solo ad hoc \\
Integridad & Sin citas inventadas; límites claros & Omite verificación \\
\bottomrule
\end{tabularx}
\end{frame}

% --- 17 Takeaways ---
\begin{frame}{Takeaways de la sesión}
\begin{enumerate}
\item \textbf{Agente $\neq$ chat:} rol, límites, herramientas y parada
\item \textbf{HITL es el estándar} en STEM académico
\item \textbf{Integridad:} ecuaciones, unidades, citas y datos se verifican
\item \textbf{Transparencia:} agent log desde el día 1
\item El CADi son \textbf{40\,h}; las \textbf{16\,h} independientes cuentan
\end{enumerate}
\end{frame}

% --- 18 TI ---
\begin{frame}{Trabajo independiente ($\sim$1\,h) --- bloque A}
Antes de la sesión 2:
\begin{itemize}
\item Lecturas / demos del \textbf{bloque A} ($\sim$1\,h; ver cronograma)
\item Revisar plantillas en \texttt{tex/plantillas/}
\item Traer bosquejo mental del \textbf{brief} de su tema EIC
\end{itemize}
Canal de dudas: \emph{(definir por edición)}.
\end{frame}

% --- 19 Prep sesión 2 ---
\begin{frame}{Prep sesión 2 --- Briefs y orquestación}
Próxima sesión (3\,h):
\begin{itemize}
\item Brief como contrato pedagógico
\item Prompt de sistema reutilizable
\item Orquestación multi-rol con puntos HITL
\end{itemize}
Traiga su \textbf{tema EIC} (aunque sea tentativo) y su nota de integridad.
\end{frame}

% --- 20 Cierre ---
\begin{frame}{Gracias}
\begin{center}
{\Large Juliho Castillo}\\[0.5em]
Escuela de Ingeniería y Ciencias\\[1em]
CADi · IA agéntica para la creación de contenido académico\\[1.5em]
{\small \emph{Verificar políticas de integridad y uso de IA con CEDDIE antes de cada edición.}}
\end{center}
\end{frame}

\end{document}
